
\documentclass[11pt,a4paper]{article}
\usepackage{Act}
\begin{document}
\input{\detokenize{/home/fenarius/Travail/Cours/cpge-info/latex/Macros.tex}}
\RDS{5}{11/04/2026}

\begin{tcolorbox}[enhanced,width=\textwidth,center upper,fontupper=\bfseries,drop shadow southwest,sharp corners]
{\sc \large Astruc} Alexandre
\end{tcolorbox}
\medskip
\begin{tabularx}{\textwidth}{p{5cm}X}
	\alertbox{\faAward}{Note}{
		\begin{itemize}[leftmargin=0pt]
			\item[\textbullet] Note : \textbf{\large 7.6}
			\item[\textbullet] Rang : \textbf{8}
			\item[\textbullet] Traité : 76 \%
		\end{itemize}
	} &
	\alertbox{\faChartLine}{Statistiques des notes}{
        \psset{xunit=1cm, yunit=1cm,fillstyle=solid}
		\begin{pspicture}(0,-0.1)(16,1.45)
		   \savedata{\data}[7.6 6.6 17.9 4.5 3.5 5.3 16.6 2.7 6.1 11.6 18.0 18.0 18.3 6.3 9.6]
		   \rput{-90}(0,0.9){\psBoxplot[barwidth=1.1cm,yunit=0.5,fillcolor=gray,linewidth=1pt]{\data}}
		   \psaxes[yAxis=false,dx=1cm,Dx=2,labelsep=1pt,linecolor=gray,xlabelFontSize=\scriptstyle](0,0)(10.1,4)
		   \psdot[dotsize=8pt,dotstyle=diamond,linecolor=black,fillstyle=solid,fillcolor=white,linewidth=1pt](3.8,0.85)
           \psdot[dotsize=6pt,dotstyle=x,linecolor=black,linewidth=3pt](5.086666666666667,0.85)
		   \end{pspicture}
	} \\
    
\end{tabularx}\\
\begin{tabularx}{\textwidth}{X}
\alertbox{\faComment}{Commentaire}
{
	Il faut impérativement faire un effort sur ton écriture notamment dans l’écriture des codes informatiques. Les exercices 3 et 4 sont à revoir en profondeur à la maison avec le corrigé. La complexité des algorithmes doit être retravaillée.
}
\end{tabularx}
\medskip
    \ding{113} \textbf{\sffamily{Résultats par thème}} \medskip \\
    \renewcommand{\arraystretch}{1.2}
    \begin{tabular}{|l|r|r|}
    \cline{2-3}
    \multicolumn{1}{l|}{} & \multicolumn{1}{|c|}{Points} & \multicolumn{1}{|c|}{Traitées} \\
    \hline
    {Diviser pour régner} & 5\% \;{\small (02/35)} & 33\% \;{\small (1/3)} \\ \hline {Comprendre un algorithme} & 100\% \;{\small (35/35)} & 100\% \;{\small (4/4)} \\ \hline {Complexité d'un algorithme} & 18\% \;{\small (16/85)} & 75\% \;{\small (6/8)} \\ \hline {Preuve de terminaison} & 0\% \;{\small (00/15)} & 0\% \;{\small (0/1)} \\ \hline {Base d'OCaml (fonctionnel)} & 56\% \;{\small (37/65)} & 83\% \;{\small (5/6)} \\ \hline {Types structurés en C et pointeurs} & 0\% \;{\small (00/10)} & 100\% \;{\small (1/1)} \\ \hline {Programmation de base en C} & 75\% \;{\small (15/20)} & 66\% \;{\small (2/3)} \\ \hline {Programmation dynamique} & 27\% \;{\small (11/40)} & 75\% \;{\small (3/4)} \\ \hline {Ordre et variant} & 36\% \;{\small (11/30)} & 100\% \;{\small (3/3)} \\ \hline \end{tabular} \\\\\medskip \\
    \ding{113} \textbf{\sffamily{Résultats par exercice}} \medskip \\
    \renewcommand{\arraystretch}{1.2}
    \begin{tabular}{|l|r|r|}
    \cline{2-3}
    \multicolumn{1}{l|}{} & \multicolumn{1}{|c|}{Points} & \multicolumn{1}{|c|}{Traitées} \\
    \hline
    Exercice {1} & 36\% \;{\small (11/30)} & 100\% \;{\small (3/3)} \\ \hline Exercice {2} & 50\% \;{\small (15/30)} & 66\% \;{\small (2/3)} \\ \hline Exercice {3} & 40\% \;{\small (47/115)} & 90\% \;{\small (10/11)} \\ \hline Exercice {4} & 33\% \;{\small (54/160)} & 62\% \;{\small (10/16)} \\ \hline \end{tabular} \\\\
    \vspace{0.5cm}\\
    \ding{113} \textbf{\sffamily{Historique des notes}} \medskip \\
    \psset{xunit=1.4cm, yunit=0.2cm}
    \begin{pspicture}(-1,-1)(12,22)

% --- Axe X : numéros de devoir ---
\psaxes[
    Dx=1,
    Dy=2,
    Ox=0,
    Oy=0,
    labels=all,
    ticks=all,
](0,0)(0,0)(7,20)


\listplot[plotstyle=line,showpoints=true,linecolor=black,linewidth=0.7pt, dotstyle=diamond*, dotsize=0.2]{%
    1 11.7
2 11.9
3 10.3
4 11.6
5 7.6
}

% Minimum de la classe
\listplot[plotstyle=line,showpoints=true,linecolor=gray,linewidth=0.7pt, dotstyle=|, dotangle=90, dotsize=0.15, linestyle=dotted]{%
    1 5.8
2 4.6
3 0.9
4 0
5 2.7
}

% Maximum de la classe
\listplot[plotstyle=line,showpoints=true,linecolor=gray,linewidth=0.7pt, dotstyle=|, dotangle=90, dotsize=0.15, linestyle=dotted]{%
    1 18.9
2 18.4
3 19.7
4 18.8
5 18.3
}

% Moyenne de la classe
\listplot[plotstyle=line,showpoints=true,linecolor=gray,linewidth=0.7pt, dotstyle=x, dotsize = 0.15, linestyle=dashed]{%
    1 12.4466666666667
2 13.1933333333333
3 10.82
4 11.5066666666667
5 10.1733333333333
}
\end{pspicture}
\pagebreak
\begin{tcolorbox}[enhanced,width=\textwidth,center upper,fontupper=\bfseries,drop shadow southwest,sharp corners]
{\sc \large Berfeuil} Rohan
\end{tcolorbox}
\medskip
\begin{tabularx}{\textwidth}{p{5cm}X}
	\alertbox{\faAward}{Note}{
		\begin{itemize}[leftmargin=0pt]
			\item[\textbullet] Note : \textbf{\large 6.6}
			\item[\textbullet] Rang : \textbf{9}
			\item[\textbullet] Traité : 48 \%
		\end{itemize}
	} &
	\alertbox{\faChartLine}{Statistiques des notes}{
        \psset{xunit=1cm, yunit=1cm,fillstyle=solid}
		\begin{pspicture}(0,-0.1)(16,1.45)
		   \savedata{\data}[7.6 6.6 17.9 4.5 3.5 5.3 16.6 2.7 6.1 11.6 18.0 18.0 18.3 6.3 9.6]
		   \rput{-90}(0,0.9){\psBoxplot[barwidth=1.1cm,yunit=0.5,fillcolor=gray,linewidth=1pt]{\data}}
		   \psaxes[yAxis=false,dx=1cm,Dx=2,labelsep=1pt,linecolor=gray,xlabelFontSize=\scriptstyle](0,0)(10.1,4)
		   \psdot[dotsize=8pt,dotstyle=diamond,linecolor=black,fillstyle=solid,fillcolor=white,linewidth=1pt](3.3,0.85)
           \psdot[dotsize=6pt,dotstyle=x,linecolor=black,linewidth=3pt](5.086666666666667,0.85)
		   \end{pspicture}
	} \\
    
\end{tabularx}\\
\begin{tabularx}{\textwidth}{X}
\alertbox{\faComment}{Commentaire}
{
	Il faut gagner en rapidité, l’exercice 4 n’a quasiment pas été traité, il faut le retravailler à la maison. Il faut retravailler la notion de complexité
}
\end{tabularx}
\medskip
    \ding{113} \textbf{\sffamily{Résultats par thème}} \medskip \\
    \renewcommand{\arraystretch}{1.2}
    \begin{tabular}{|l|r|r|}
    \cline{2-3}
    \multicolumn{1}{l|}{} & \multicolumn{1}{|c|}{Points} & \multicolumn{1}{|c|}{Traitées} \\
    \hline
    {Diviser pour régner} & 0\% \;{\small (00/35)} & 0\% \;{\small (0/3)} \\ \hline {Comprendre un algorithme} & 100\% \;{\small (35/35)} & 100\% \;{\small (4/4)} \\ \hline {Complexité d'un algorithme} & 11\% \;{\small (10/85)} & 37\% \;{\small (3/8)} \\ \hline {Preuve de terminaison} & 0\% \;{\small (00/15)} & 0\% \;{\small (0/1)} \\ \hline {Base d'OCaml (fonctionnel)} & 53\% \;{\small (35/65)} & 83\% \;{\small (5/6)} \\ \hline {Types structurés en C et pointeurs} & 0\% \;{\small (00/10)} & 0\% \;{\small (0/1)} \\ \hline {Programmation de base en C} & 75\% \;{\small (15/20)} & 66\% \;{\small (2/3)} \\ \hline {Programmation dynamique} & 0\% \;{\small (00/40)} & 0\% \;{\small (0/4)} \\ \hline {Ordre et variant} & 50\% \;{\small (15/30)} & 66\% \;{\small (2/3)} \\ \hline \end{tabular} \\\\\medskip \\
    \ding{113} \textbf{\sffamily{Résultats par exercice}} \medskip \\
    \renewcommand{\arraystretch}{1.2}
    \begin{tabular}{|l|r|r|}
    \cline{2-3}
    \multicolumn{1}{l|}{} & \multicolumn{1}{|c|}{Points} & \multicolumn{1}{|c|}{Traitées} \\
    \hline
    Exercice {1} & 50\% \;{\small (15/30)} & 66\% \;{\small (2/3)} \\ \hline Exercice {2} & 50\% \;{\small (15/30)} & 66\% \;{\small (2/3)} \\ \hline Exercice {3} & 39\% \;{\small (45/115)} & 72\% \;{\small (8/11)} \\ \hline Exercice {4} & 21\% \;{\small (35/160)} & 25\% \;{\small (4/16)} \\ \hline \end{tabular} \\\\
    \vspace{0.5cm}\\
    \ding{113} \textbf{\sffamily{Historique des notes}} \medskip \\
    \psset{xunit=1.4cm, yunit=0.2cm}
    \begin{pspicture}(-1,-1)(12,22)

% --- Axe X : numéros de devoir ---
\psaxes[
    Dx=1,
    Dy=2,
    Ox=0,
    Oy=0,
    labels=all,
    ticks=all,
](0,0)(0,0)(7,20)


\listplot[plotstyle=line,showpoints=true,linecolor=black,linewidth=0.7pt, dotstyle=diamond*, dotsize=0.2]{%
    1 11.2
2 15.4
3 16.4
4 14.5
5 6.6
}

% Minimum de la classe
\listplot[plotstyle=line,showpoints=true,linecolor=gray,linewidth=0.7pt, dotstyle=|, dotangle=90, dotsize=0.15, linestyle=dotted]{%
    1 5.8
2 4.6
3 0.9
4 0
5 2.7
}

% Maximum de la classe
\listplot[plotstyle=line,showpoints=true,linecolor=gray,linewidth=0.7pt, dotstyle=|, dotangle=90, dotsize=0.15, linestyle=dotted]{%
    1 18.9
2 18.4
3 19.7
4 18.8
5 18.3
}

% Moyenne de la classe
\listplot[plotstyle=line,showpoints=true,linecolor=gray,linewidth=0.7pt, dotstyle=x, dotsize = 0.15, linestyle=dashed]{%
    1 12.4466666666667
2 13.1933333333333
3 10.82
4 11.5066666666667
5 10.1733333333333
}
\end{pspicture}
\pagebreak
\begin{tcolorbox}[enhanced,width=\textwidth,center upper,fontupper=\bfseries,drop shadow southwest,sharp corners]
{\sc \large Body} Timothée
\end{tcolorbox}
\medskip
\begin{tabularx}{\textwidth}{p{5cm}X}
	\alertbox{\faAward}{Note}{
		\begin{itemize}[leftmargin=0pt]
			\item[\textbullet] Note : \textbf{\large 17.9}
			\item[\textbullet] Rang : \textbf{4}
			\item[\textbullet] Traité : 100 \%
		\end{itemize}
	} &
	\alertbox{\faChartLine}{Statistiques des notes}{
        \psset{xunit=1cm, yunit=1cm,fillstyle=solid}
		\begin{pspicture}(0,-0.1)(16,1.45)
		   \savedata{\data}[7.6 6.6 17.9 4.5 3.5 5.3 16.6 2.7 6.1 11.6 18.0 18.0 18.3 6.3 9.6]
		   \rput{-90}(0,0.9){\psBoxplot[barwidth=1.1cm,yunit=0.5,fillcolor=gray,linewidth=1pt]{\data}}
		   \psaxes[yAxis=false,dx=1cm,Dx=2,labelsep=1pt,linecolor=gray,xlabelFontSize=\scriptstyle](0,0)(10.1,4)
		   \psdot[dotsize=8pt,dotstyle=diamond,linecolor=black,fillstyle=solid,fillcolor=white,linewidth=1pt](8.95,0.85)
           \psdot[dotsize=6pt,dotstyle=x,linecolor=black,linewidth=3pt](5.086666666666667,0.85)
		   \end{pspicture}
	} \\
    
\end{tabularx}\\
\begin{tabularx}{\textwidth}{X}
\alertbox{\faComment}{Commentaire}
{
	Très bon travail, certains calcul de complexité doivent être plus détaillés. Fais attention si un algorithme appelle une fonction f puis une fonction g sa complexité est la somme de celle de f et de g (et pas le produit) !
}
\end{tabularx}
\medskip
    \ding{113} \textbf{\sffamily{Résultats par thème}} \medskip \\
    \renewcommand{\arraystretch}{1.2}
    \begin{tabular}{|l|r|r|}
    \cline{2-3}
    \multicolumn{1}{l|}{} & \multicolumn{1}{|c|}{Points} & \multicolumn{1}{|c|}{Traitées} \\
    \hline
    {Diviser pour régner} & 71\% \;{\small (25/35)} & 100\% \;{\small (3/3)} \\ \hline {Comprendre un algorithme} & 100\% \;{\small (35/35)} & 100\% \;{\small (4/4)} \\ \hline {Complexité d'un algorithme} & 72\% \;{\small (62/85)} & 100\% \;{\small (8/8)} \\ \hline {Preuve de terminaison} & 100\% \;{\small (15/15)} & 100\% \;{\small (1/1)} \\ \hline {Base d'OCaml (fonctionnel)} & 100\% \;{\small (65/65)} & 100\% \;{\small (6/6)} \\ \hline {Types structurés en C et pointeurs} & 100\% \;{\small (10/10)} & 100\% \;{\small (1/1)} \\ \hline {Programmation de base en C} & 100\% \;{\small (20/20)} & 100\% \;{\small (3/3)} \\ \hline {Programmation dynamique} & 100\% \;{\small (40/40)} & 100\% \;{\small (4/4)} \\ \hline {Ordre et variant} & 90\% \;{\small (27/30)} & 100\% \;{\small (3/3)} \\ \hline \end{tabular} \\\\\medskip \\
    \ding{113} \textbf{\sffamily{Résultats par exercice}} \medskip \\
    \renewcommand{\arraystretch}{1.2}
    \begin{tabular}{|l|r|r|}
    \cline{2-3}
    \multicolumn{1}{l|}{} & \multicolumn{1}{|c|}{Points} & \multicolumn{1}{|c|}{Traitées} \\
    \hline
    Exercice {1} & 90\% \;{\small (27/30)} & 100\% \;{\small (3/3)} \\ \hline Exercice {2} & 100\% \;{\small (30/30)} & 100\% \;{\small (3/3)} \\ \hline Exercice {3} & 87\% \;{\small (101/115)} & 100\% \;{\small (11/11)} \\ \hline Exercice {4} & 88\% \;{\small (141/160)} & 100\% \;{\small (16/16)} \\ \hline \end{tabular} \\\\
    \vspace{0.5cm}\\
    \ding{113} \textbf{\sffamily{Historique des notes}} \medskip \\
    \psset{xunit=1.4cm, yunit=0.2cm}
    \begin{pspicture}(-1,-1)(12,22)

% --- Axe X : numéros de devoir ---
\psaxes[
    Dx=1,
    Dy=2,
    Ox=0,
    Oy=0,
    labels=all,
    ticks=all,
](0,0)(0,0)(7,20)


\listplot[plotstyle=line,showpoints=true,linecolor=black,linewidth=0.7pt, dotstyle=diamond*, dotsize=0.2]{%
    1 18.9
2 18.4
3 19.2
4 18.8
5 17.9
}

% Minimum de la classe
\listplot[plotstyle=line,showpoints=true,linecolor=gray,linewidth=0.7pt, dotstyle=|, dotangle=90, dotsize=0.15, linestyle=dotted]{%
    1 5.8
2 4.6
3 0.9
4 0
5 2.7
}

% Maximum de la classe
\listplot[plotstyle=line,showpoints=true,linecolor=gray,linewidth=0.7pt, dotstyle=|, dotangle=90, dotsize=0.15, linestyle=dotted]{%
    1 18.9
2 18.4
3 19.7
4 18.8
5 18.3
}

% Moyenne de la classe
\listplot[plotstyle=line,showpoints=true,linecolor=gray,linewidth=0.7pt, dotstyle=x, dotsize = 0.15, linestyle=dashed]{%
    1 12.4466666666667
2 13.1933333333333
3 10.82
4 11.5066666666667
5 10.1733333333333
}
\end{pspicture}
\pagebreak
\begin{tcolorbox}[enhanced,width=\textwidth,center upper,fontupper=\bfseries,drop shadow southwest,sharp corners]
{\sc \large Boucher} Mathis
\end{tcolorbox}
\medskip
\begin{tabularx}{\textwidth}{p{5cm}X}
	\alertbox{\faAward}{Note}{
		\begin{itemize}[leftmargin=0pt]
			\item[\textbullet] Note : \textbf{\large 4.5}
			\item[\textbullet] Rang : \textbf{13}
			\item[\textbullet] Traité : 55 \%
		\end{itemize}
	} &
	\alertbox{\faChartLine}{Statistiques des notes}{
        \psset{xunit=1cm, yunit=1cm,fillstyle=solid}
		\begin{pspicture}(0,-0.1)(16,1.45)
		   \savedata{\data}[7.6 6.6 17.9 4.5 3.5 5.3 16.6 2.7 6.1 11.6 18.0 18.0 18.3 6.3 9.6]
		   \rput{-90}(0,0.9){\psBoxplot[barwidth=1.1cm,yunit=0.5,fillcolor=gray,linewidth=1pt]{\data}}
		   \psaxes[yAxis=false,dx=1cm,Dx=2,labelsep=1pt,linecolor=gray,xlabelFontSize=\scriptstyle](0,0)(10.1,4)
		   \psdot[dotsize=8pt,dotstyle=diamond,linecolor=black,fillstyle=solid,fillcolor=white,linewidth=1pt](2.25,0.85)
           \psdot[dotsize=6pt,dotstyle=x,linecolor=black,linewidth=3pt](5.086666666666667,0.85)
		   \end{pspicture}
	} \\
    
\end{tabularx}\\
\begin{tabularx}{\textwidth}{X}
\alertbox{\faComment}{Commentaire}
{
	Il faut gagner en rapidité l’exercice 4 n’a presque pas été traité. La structure de tas n’est pas acquise.
}
\end{tabularx}
\medskip
    \ding{113} \textbf{\sffamily{Résultats par thème}} \medskip \\
    \renewcommand{\arraystretch}{1.2}
    \begin{tabular}{|l|r|r|}
    \cline{2-3}
    \multicolumn{1}{l|}{} & \multicolumn{1}{|c|}{Points} & \multicolumn{1}{|c|}{Traitées} \\
    \hline
    {Diviser pour régner} & 0\% \;{\small (00/35)} & 0\% \;{\small (0/3)} \\ \hline {Comprendre un algorithme} & 25\% \;{\small (09/35)} & 75\% \;{\small (3/4)} \\ \hline {Complexité d'un algorithme} & 16\% \;{\small (14/85)} & 62\% \;{\small (5/8)} \\ \hline {Preuve de terminaison} & 60\% \;{\small (09/15)} & 100\% \;{\small (1/1)} \\ \hline {Base d'OCaml (fonctionnel)} & 40\% \;{\small (26/65)} & 83\% \;{\small (5/6)} \\ \hline {Types structurés en C et pointeurs} & 0\% \;{\small (00/10)} & 0\% \;{\small (0/1)} \\ \hline {Programmation de base en C} & 25\% \;{\small (05/20)} & 33\% \;{\small (1/3)} \\ \hline {Programmation dynamique} & 0\% \;{\small (00/40)} & 0\% \;{\small (0/4)} \\ \hline {Ordre et variant} & 43\% \;{\small (13/30)} & 100\% \;{\small (3/3)} \\ \hline \end{tabular} \\\\\medskip \\
    \ding{113} \textbf{\sffamily{Résultats par exercice}} \medskip \\
    \renewcommand{\arraystretch}{1.2}
    \begin{tabular}{|l|r|r|}
    \cline{2-3}
    \multicolumn{1}{l|}{} & \multicolumn{1}{|c|}{Points} & \multicolumn{1}{|c|}{Traitées} \\
    \hline
    Exercice {1} & 43\% \;{\small (13/30)} & 100\% \;{\small (3/3)} \\ \hline Exercice {2} & 80\% \;{\small (24/30)} & 100\% \;{\small (3/3)} \\ \hline Exercice {3} & 26\% \;{\small (30/115)} & 90\% \;{\small (10/11)} \\ \hline Exercice {4} & 5\% \;{\small (09/160)} & 12\% \;{\small (2/16)} \\ \hline \end{tabular} \\\\
    \vspace{0.5cm}\\
    \ding{113} \textbf{\sffamily{Historique des notes}} \medskip \\
    \psset{xunit=1.4cm, yunit=0.2cm}
    \begin{pspicture}(-1,-1)(12,22)

% --- Axe X : numéros de devoir ---
\psaxes[
    Dx=1,
    Dy=2,
    Ox=0,
    Oy=0,
    labels=all,
    ticks=all,
](0,0)(0,0)(7,20)


\listplot[plotstyle=line,showpoints=true,linecolor=black,linewidth=0.7pt, dotstyle=diamond*, dotsize=0.2]{%
    1 8.4
2 11.8
3 8
4 8.4
5 4.5
}

% Minimum de la classe
\listplot[plotstyle=line,showpoints=true,linecolor=gray,linewidth=0.7pt, dotstyle=|, dotangle=90, dotsize=0.15, linestyle=dotted]{%
    1 5.8
2 4.6
3 0.9
4 0
5 2.7
}

% Maximum de la classe
\listplot[plotstyle=line,showpoints=true,linecolor=gray,linewidth=0.7pt, dotstyle=|, dotangle=90, dotsize=0.15, linestyle=dotted]{%
    1 18.9
2 18.4
3 19.7
4 18.8
5 18.3
}

% Moyenne de la classe
\listplot[plotstyle=line,showpoints=true,linecolor=gray,linewidth=0.7pt, dotstyle=x, dotsize = 0.15, linestyle=dashed]{%
    1 12.4466666666667
2 13.1933333333333
3 10.82
4 11.5066666666667
5 10.1733333333333
}
\end{pspicture}
\pagebreak
\begin{tcolorbox}[enhanced,width=\textwidth,center upper,fontupper=\bfseries,drop shadow southwest,sharp corners]
{\sc \large Chane-lock} Maxime
\end{tcolorbox}
\medskip
\begin{tabularx}{\textwidth}{p{5cm}X}
	\alertbox{\faAward}{Note}{
		\begin{itemize}[leftmargin=0pt]
			\item[\textbullet] Note : \textbf{\large 3.5}
			\item[\textbullet] Rang : \textbf{14}
			\item[\textbullet] Traité : 45 \%
		\end{itemize}
	} &
	\alertbox{\faChartLine}{Statistiques des notes}{
        \psset{xunit=1cm, yunit=1cm,fillstyle=solid}
		\begin{pspicture}(0,-0.1)(16,1.45)
		   \savedata{\data}[7.6 6.6 17.9 4.5 3.5 5.3 16.6 2.7 6.1 11.6 18.0 18.0 18.3 6.3 9.6]
		   \rput{-90}(0,0.9){\psBoxplot[barwidth=1.1cm,yunit=0.5,fillcolor=gray,linewidth=1pt]{\data}}
		   \psaxes[yAxis=false,dx=1cm,Dx=2,labelsep=1pt,linecolor=gray,xlabelFontSize=\scriptstyle](0,0)(10.1,4)
		   \psdot[dotsize=8pt,dotstyle=diamond,linecolor=black,fillstyle=solid,fillcolor=white,linewidth=1pt](1.75,0.85)
           \psdot[dotsize=6pt,dotstyle=x,linecolor=black,linewidth=3pt](5.086666666666667,0.85)
		   \end{pspicture}
	} \\
    
\end{tabularx}\\
\begin{tabularx}{\textwidth}{X}
\alertbox{\faComment}{Commentaire}
{
	Résultats inquiétants, il faut revoir le cours et reprendre à la maison (avec le corrigé) les exercices 3 et 4.
}
\end{tabularx}
\medskip
    \ding{113} \textbf{\sffamily{Résultats par thème}} \medskip \\
    \renewcommand{\arraystretch}{1.2}
    \begin{tabular}{|l|r|r|}
    \cline{2-3}
    \multicolumn{1}{l|}{} & \multicolumn{1}{|c|}{Points} & \multicolumn{1}{|c|}{Traitées} \\
    \hline
    {Diviser pour régner} & 0\% \;{\small (00/35)} & 66\% \;{\small (2/3)} \\ \hline {Comprendre un algorithme} & 11\% \;{\small (04/35)} & 75\% \;{\small (3/4)} \\ \hline {Complexité d'un algorithme} & 0\% \;{\small (00/85)} & 0\% \;{\small (0/8)} \\ \hline {Preuve de terminaison} & 20\% \;{\small (03/15)} & 100\% \;{\small (1/1)} \\ \hline {Base d'OCaml (fonctionnel)} & 36\% \;{\small (24/65)} & 66\% \;{\small (4/6)} \\ \hline {Types structurés en C et pointeurs} & 0\% \;{\small (00/10)} & 0\% \;{\small (0/1)} \\ \hline {Programmation de base en C} & 65\% \;{\small (13/20)} & 66\% \;{\small (2/3)} \\ \hline {Programmation dynamique} & 0\% \;{\small (00/40)} & 0\% \;{\small (0/4)} \\ \hline {Ordre et variant} & 46\% \;{\small (14/30)} & 100\% \;{\small (3/3)} \\ \hline \end{tabular} \\\\\medskip \\
    \ding{113} \textbf{\sffamily{Résultats par exercice}} \medskip \\
    \renewcommand{\arraystretch}{1.2}
    \begin{tabular}{|l|r|r|}
    \cline{2-3}
    \multicolumn{1}{l|}{} & \multicolumn{1}{|c|}{Points} & \multicolumn{1}{|c|}{Traitées} \\
    \hline
    Exercice {1} & 46\% \;{\small (14/30)} & 100\% \;{\small (3/3)} \\ \hline Exercice {2} & 43\% \;{\small (13/30)} & 100\% \;{\small (3/3)} \\ \hline Exercice {3} & 15\% \;{\small (18/115)} & 36\% \;{\small (4/11)} \\ \hline Exercice {4} & 8\% \;{\small (13/160)} & 31\% \;{\small (5/16)} \\ \hline \end{tabular} \\\\
    \vspace{0.5cm}\\
    \ding{113} \textbf{\sffamily{Historique des notes}} \medskip \\
    \psset{xunit=1.4cm, yunit=0.2cm}
    \begin{pspicture}(-1,-1)(12,22)

% --- Axe X : numéros de devoir ---
\psaxes[
    Dx=1,
    Dy=2,
    Ox=0,
    Oy=0,
    labels=all,
    ticks=all,
](0,0)(0,0)(7,20)


\listplot[plotstyle=line,showpoints=true,linecolor=black,linewidth=0.7pt, dotstyle=diamond*, dotsize=0.2]{%
    1 8.4
2 9
3 8.1
4 8.2
5 3.5
}

% Minimum de la classe
\listplot[plotstyle=line,showpoints=true,linecolor=gray,linewidth=0.7pt, dotstyle=|, dotangle=90, dotsize=0.15, linestyle=dotted]{%
    1 5.8
2 4.6
3 0.9
4 0
5 2.7
}

% Maximum de la classe
\listplot[plotstyle=line,showpoints=true,linecolor=gray,linewidth=0.7pt, dotstyle=|, dotangle=90, dotsize=0.15, linestyle=dotted]{%
    1 18.9
2 18.4
3 19.7
4 18.8
5 18.3
}

% Moyenne de la classe
\listplot[plotstyle=line,showpoints=true,linecolor=gray,linewidth=0.7pt, dotstyle=x, dotsize = 0.15, linestyle=dashed]{%
    1 12.4466666666667
2 13.1933333333333
3 10.82
4 11.5066666666667
5 10.1733333333333
}
\end{pspicture}
\pagebreak
\begin{tcolorbox}[enhanced,width=\textwidth,center upper,fontupper=\bfseries,drop shadow southwest,sharp corners]
{\sc \large Courounadin-Mouny} Maxence
\end{tcolorbox}
\medskip
\begin{tabularx}{\textwidth}{p{5cm}X}
	\alertbox{\faAward}{Note}{
		\begin{itemize}[leftmargin=0pt]
			\item[\textbullet] Note : \textbf{\large 5.3}
			\item[\textbullet] Rang : \textbf{12}
			\item[\textbullet] Traité : 76 \%
		\end{itemize}
	} &
	\alertbox{\faChartLine}{Statistiques des notes}{
        \psset{xunit=1cm, yunit=1cm,fillstyle=solid}
		\begin{pspicture}(0,-0.1)(16,1.45)
		   \savedata{\data}[7.6 6.6 17.9 4.5 3.5 5.3 16.6 2.7 6.1 11.6 18.0 18.0 18.3 6.3 9.6]
		   \rput{-90}(0,0.9){\psBoxplot[barwidth=1.1cm,yunit=0.5,fillcolor=gray,linewidth=1pt]{\data}}
		   \psaxes[yAxis=false,dx=1cm,Dx=2,labelsep=1pt,linecolor=gray,xlabelFontSize=\scriptstyle](0,0)(10.1,4)
		   \psdot[dotsize=8pt,dotstyle=diamond,linecolor=black,fillstyle=solid,fillcolor=white,linewidth=1pt](2.65,0.85)
           \psdot[dotsize=6pt,dotstyle=x,linecolor=black,linewidth=3pt](5.086666666666667,0.85)
		   \end{pspicture}
	} \\
    
\end{tabularx}\\
\begin{tabularx}{\textwidth}{X}
\alertbox{\faComment}{Commentaire}
{
	Il faut revoir le cours sur les relations d’ordre. Reprends à la maison les exercices 3 et 4 en codant et en testant les fonctions demandées. Attention à la confusion entre tableaux et listes en Ocaml.
}
\end{tabularx}
\medskip
    \ding{113} \textbf{\sffamily{Résultats par thème}} \medskip \\
    \renewcommand{\arraystretch}{1.2}
    \begin{tabular}{|l|r|r|}
    \cline{2-3}
    \multicolumn{1}{l|}{} & \multicolumn{1}{|c|}{Points} & \multicolumn{1}{|c|}{Traitées} \\
    \hline
    {Diviser pour régner} & 0\% \;{\small (00/35)} & 66\% \;{\small (2/3)} \\ \hline {Comprendre un algorithme} & 42\% \;{\small (15/35)} & 100\% \;{\small (4/4)} \\ \hline {Complexité d'un algorithme} & 28\% \;{\small (24/85)} & 87\% \;{\small (7/8)} \\ \hline {Preuve de terminaison} & 0\% \;{\small (00/15)} & 100\% \;{\small (1/1)} \\ \hline {Base d'OCaml (fonctionnel)} & 27\% \;{\small (18/65)} & 83\% \;{\small (5/6)} \\ \hline {Types structurés en C et pointeurs} & 100\% \;{\small (10/10)} & 100\% \;{\small (1/1)} \\ \hline {Programmation de base en C} & 75\% \;{\small (15/20)} & 66\% \;{\small (2/3)} \\ \hline {Programmation dynamique} & 0\% \;{\small (00/40)} & 0\% \;{\small (0/4)} \\ \hline {Ordre et variant} & 20\% \;{\small (06/30)} & 100\% \;{\small (3/3)} \\ \hline \end{tabular} \\\\\medskip \\
    \ding{113} \textbf{\sffamily{Résultats par exercice}} \medskip \\
    \renewcommand{\arraystretch}{1.2}
    \begin{tabular}{|l|r|r|}
    \cline{2-3}
    \multicolumn{1}{l|}{} & \multicolumn{1}{|c|}{Points} & \multicolumn{1}{|c|}{Traitées} \\
    \hline
    Exercice {1} & 20\% \;{\small (06/30)} & 100\% \;{\small (3/3)} \\ \hline Exercice {2} & 50\% \;{\small (15/30)} & 100\% \;{\small (3/3)} \\ \hline Exercice {3} & 19\% \;{\small (22/115)} & 90\% \;{\small (10/11)} \\ \hline Exercice {4} & 28\% \;{\small (45/160)} & 56\% \;{\small (9/16)} \\ \hline \end{tabular} \\\\
    \vspace{0.5cm}\\
    \ding{113} \textbf{\sffamily{Historique des notes}} \medskip \\
    \psset{xunit=1.4cm, yunit=0.2cm}
    \begin{pspicture}(-1,-1)(12,22)

% --- Axe X : numéros de devoir ---
\psaxes[
    Dx=1,
    Dy=2,
    Ox=0,
    Oy=0,
    labels=all,
    ticks=all,
](0,0)(0,0)(7,20)


\listplot[plotstyle=line,showpoints=true,linecolor=black,linewidth=0.7pt, dotstyle=diamond*, dotsize=0.2]{%
    1 10.9
2 11.1
3 4.8
4 7.6
5 5.3
}

% Minimum de la classe
\listplot[plotstyle=line,showpoints=true,linecolor=gray,linewidth=0.7pt, dotstyle=|, dotangle=90, dotsize=0.15, linestyle=dotted]{%
    1 5.8
2 4.6
3 0.9
4 0
5 2.7
}

% Maximum de la classe
\listplot[plotstyle=line,showpoints=true,linecolor=gray,linewidth=0.7pt, dotstyle=|, dotangle=90, dotsize=0.15, linestyle=dotted]{%
    1 18.9
2 18.4
3 19.7
4 18.8
5 18.3
}

% Moyenne de la classe
\listplot[plotstyle=line,showpoints=true,linecolor=gray,linewidth=0.7pt, dotstyle=x, dotsize = 0.15, linestyle=dashed]{%
    1 12.4466666666667
2 13.1933333333333
3 10.82
4 11.5066666666667
5 10.1733333333333
}
\end{pspicture}
\pagebreak
\begin{tcolorbox}[enhanced,width=\textwidth,center upper,fontupper=\bfseries,drop shadow southwest,sharp corners]
{\sc \large Dominguez} Raphaël
\end{tcolorbox}
\medskip
\begin{tabularx}{\textwidth}{p{5cm}X}
	\alertbox{\faAward}{Note}{
		\begin{itemize}[leftmargin=0pt]
			\item[\textbullet] Note : \textbf{\large 16.6}
			\item[\textbullet] Rang : \textbf{5}
			\item[\textbullet] Traité : 91 \%
		\end{itemize}
	} &
	\alertbox{\faChartLine}{Statistiques des notes}{
        \psset{xunit=1cm, yunit=1cm,fillstyle=solid}
		\begin{pspicture}(0,-0.1)(16,1.45)
		   \savedata{\data}[7.6 6.6 17.9 4.5 3.5 5.3 16.6 2.7 6.1 11.6 18.0 18.0 18.3 6.3 9.6]
		   \rput{-90}(0,0.9){\psBoxplot[barwidth=1.1cm,yunit=0.5,fillcolor=gray,linewidth=1pt]{\data}}
		   \psaxes[yAxis=false,dx=1cm,Dx=2,labelsep=1pt,linecolor=gray,xlabelFontSize=\scriptstyle](0,0)(10.1,4)
		   \psdot[dotsize=8pt,dotstyle=diamond,linecolor=black,fillstyle=solid,fillcolor=white,linewidth=1pt](8.3,0.85)
           \psdot[dotsize=6pt,dotstyle=x,linecolor=black,linewidth=3pt](5.086666666666667,0.85)
		   \end{pspicture}
	} \\
    
\end{tabularx}\\
\begin{tabularx}{\textwidth}{X}
\alertbox{\faComment}{Commentaire}
{
	Très bon travail. Attention à la structure de tas binaire qui ne semble pas maitrisée (revoir les algorithmes d’insertion et d’extraction)
}
\end{tabularx}
\medskip
    \ding{113} \textbf{\sffamily{Résultats par thème}} \medskip \\
    \renewcommand{\arraystretch}{1.2}
    \begin{tabular}{|l|r|r|}
    \cline{2-3}
    \multicolumn{1}{l|}{} & \multicolumn{1}{|c|}{Points} & \multicolumn{1}{|c|}{Traitées} \\
    \hline
    {Diviser pour régner} & 100\% \;{\small (35/35)} & 100\% \;{\small (3/3)} \\ \hline {Comprendre un algorithme} & 71\% \;{\small (25/35)} & 100\% \;{\small (4/4)} \\ \hline {Complexité d'un algorithme} & 70\% \;{\small (60/85)} & 75\% \;{\small (6/8)} \\ \hline {Preuve de terminaison} & 80\% \;{\small (12/15)} & 100\% \;{\small (1/1)} \\ \hline {Base d'OCaml (fonctionnel)} & 95\% \;{\small (62/65)} & 100\% \;{\small (6/6)} \\ \hline {Types structurés en C et pointeurs} & 100\% \;{\small (10/10)} & 100\% \;{\small (1/1)} \\ \hline {Programmation de base en C} & 100\% \;{\small (20/20)} & 100\% \;{\small (3/3)} \\ \hline {Programmation dynamique} & 60\% \;{\small (24/40)} & 75\% \;{\small (3/4)} \\ \hline {Ordre et variant} & 100\% \;{\small (30/30)} & 100\% \;{\small (3/3)} \\ \hline \end{tabular} \\\\\medskip \\
    \ding{113} \textbf{\sffamily{Résultats par exercice}} \medskip \\
    \renewcommand{\arraystretch}{1.2}
    \begin{tabular}{|l|r|r|}
    \cline{2-3}
    \multicolumn{1}{l|}{} & \multicolumn{1}{|c|}{Points} & \multicolumn{1}{|c|}{Traitées} \\
    \hline
    Exercice {1} & 100\% \;{\small (30/30)} & 100\% \;{\small (3/3)} \\ \hline Exercice {2} & 90\% \;{\small (27/30)} & 100\% \;{\small (3/3)} \\ \hline Exercice {3} & 88\% \;{\small (102/115)} & 100\% \;{\small (11/11)} \\ \hline Exercice {4} & 74\% \;{\small (119/160)} & 81\% \;{\small (13/16)} \\ \hline \end{tabular} \\\\
    \vspace{0.5cm}\\
    \ding{113} \textbf{\sffamily{Historique des notes}} \medskip \\
    \psset{xunit=1.4cm, yunit=0.2cm}
    \begin{pspicture}(-1,-1)(12,22)

% --- Axe X : numéros de devoir ---
\psaxes[
    Dx=1,
    Dy=2,
    Ox=0,
    Oy=0,
    labels=all,
    ticks=all,
](0,0)(0,0)(7,20)


\listplot[plotstyle=line,showpoints=true,linecolor=black,linewidth=0.7pt, dotstyle=diamond*, dotsize=0.2]{%
    1 15.7
2 16.1
3 6.6
4 14.1
5 16.6
}

% Minimum de la classe
\listplot[plotstyle=line,showpoints=true,linecolor=gray,linewidth=0.7pt, dotstyle=|, dotangle=90, dotsize=0.15, linestyle=dotted]{%
    1 5.8
2 4.6
3 0.9
4 0
5 2.7
}

% Maximum de la classe
\listplot[plotstyle=line,showpoints=true,linecolor=gray,linewidth=0.7pt, dotstyle=|, dotangle=90, dotsize=0.15, linestyle=dotted]{%
    1 18.9
2 18.4
3 19.7
4 18.8
5 18.3
}

% Moyenne de la classe
\listplot[plotstyle=line,showpoints=true,linecolor=gray,linewidth=0.7pt, dotstyle=x, dotsize = 0.15, linestyle=dashed]{%
    1 12.4466666666667
2 13.1933333333333
3 10.82
4 11.5066666666667
5 10.1733333333333
}
\end{pspicture}
\pagebreak
\begin{tcolorbox}[enhanced,width=\textwidth,center upper,fontupper=\bfseries,drop shadow southwest,sharp corners]
{\sc \large Garbal} Alizée
\end{tcolorbox}
\medskip
\begin{tabularx}{\textwidth}{p{5cm}X}
	\alertbox{\faAward}{Note}{
		\begin{itemize}[leftmargin=0pt]
			\item[\textbullet] Note : \textbf{\large 2.7}
			\item[\textbullet] Rang : \textbf{15}
			\item[\textbullet] Traité : 88 \%
		\end{itemize}
	} &
	\alertbox{\faChartLine}{Statistiques des notes}{
        \psset{xunit=1cm, yunit=1cm,fillstyle=solid}
		\begin{pspicture}(0,-0.1)(16,1.45)
		   \savedata{\data}[7.6 6.6 17.9 4.5 3.5 5.3 16.6 2.7 6.1 11.6 18.0 18.0 18.3 6.3 9.6]
		   \rput{-90}(0,0.9){\psBoxplot[barwidth=1.1cm,yunit=0.5,fillcolor=gray,linewidth=1pt]{\data}}
		   \psaxes[yAxis=false,dx=1cm,Dx=2,labelsep=1pt,linecolor=gray,xlabelFontSize=\scriptstyle](0,0)(10.1,4)
		   \psdot[dotsize=8pt,dotstyle=diamond,linecolor=black,fillstyle=solid,fillcolor=white,linewidth=1pt](1.35,0.85)
           \psdot[dotsize=6pt,dotstyle=x,linecolor=black,linewidth=3pt](5.086666666666667,0.85)
		   \end{pspicture}
	} \\
    
\end{tabularx}\\
\begin{tabularx}{\textwidth}{X}
\alertbox{\faComment}{Commentaire}
{
	Le cours a été appris mais tu as du mal à restituer les raisonnements dans les preuves. L’exerice 3 est à retravailler à la maison, Dans l’exercice 4, il faut relire l’énoncé et bien comprendre les définitions de saut et de valeur d’un saut.
}
\end{tabularx}
\medskip
    \ding{113} \textbf{\sffamily{Résultats par thème}} \medskip \\
    \renewcommand{\arraystretch}{1.2}
    \begin{tabular}{|l|r|r|}
    \cline{2-3}
    \multicolumn{1}{l|}{} & \multicolumn{1}{|c|}{Points} & \multicolumn{1}{|c|}{Traitées} \\
    \hline
    {Diviser pour régner} & 17\% \;{\small (06/35)} & 100\% \;{\small (3/3)} \\ \hline {Comprendre un algorithme} & 14\% \;{\small (05/35)} & 75\% \;{\small (3/4)} \\ \hline {Complexité d'un algorithme} & 0\% \;{\small (00/85)} & 75\% \;{\small (6/8)} \\ \hline {Preuve de terminaison} & 0\% \;{\small (00/15)} & 100\% \;{\small (1/1)} \\ \hline {Base d'OCaml (fonctionnel)} & 21\% \;{\small (14/65)} & 100\% \;{\small (6/6)} \\ \hline {Types structurés en C et pointeurs} & 0\% \;{\small (00/10)} & 0\% \;{\small (0/1)} \\ \hline {Programmation de base en C} & 50\% \;{\small (10/20)} & 100\% \;{\small (3/3)} \\ \hline {Programmation dynamique} & 0\% \;{\small (00/40)} & 100\% \;{\small (4/4)} \\ \hline {Ordre et variant} & 33\% \;{\small (10/30)} & 100\% \;{\small (3/3)} \\ \hline \end{tabular} \\\\\medskip \\
    \ding{113} \textbf{\sffamily{Résultats par exercice}} \medskip \\
    \renewcommand{\arraystretch}{1.2}
    \begin{tabular}{|l|r|r|}
    \cline{2-3}
    \multicolumn{1}{l|}{} & \multicolumn{1}{|c|}{Points} & \multicolumn{1}{|c|}{Traitées} \\
    \hline
    Exercice {1} & 33\% \;{\small (10/30)} & 100\% \;{\small (3/3)} \\ \hline Exercice {2} & 30\% \;{\small (09/30)} & 100\% \;{\small (3/3)} \\ \hline Exercice {3} & 5\% \;{\small (06/115)} & 90\% \;{\small (10/11)} \\ \hline Exercice {4} & 12\% \;{\small (20/160)} & 81\% \;{\small (13/16)} \\ \hline \end{tabular} \\\\
    \vspace{0.5cm}\\
    \ding{113} \textbf{\sffamily{Historique des notes}} \medskip \\
    \psset{xunit=1.4cm, yunit=0.2cm}
    \begin{pspicture}(-1,-1)(12,22)

% --- Axe X : numéros de devoir ---
\psaxes[
    Dx=1,
    Dy=2,
    Ox=0,
    Oy=0,
    labels=all,
    ticks=all,
](0,0)(0,0)(7,20)


\listplot[plotstyle=line,showpoints=true,linecolor=black,linewidth=0.7pt, dotstyle=diamond*, dotsize=0.2]{%
    1 5.8
2 4.6
3 0.9
4 8.1
5 2.7
}

% Minimum de la classe
\listplot[plotstyle=line,showpoints=true,linecolor=gray,linewidth=0.7pt, dotstyle=|, dotangle=90, dotsize=0.15, linestyle=dotted]{%
    1 5.8
2 4.6
3 0.9
4 0
5 2.7
}

% Maximum de la classe
\listplot[plotstyle=line,showpoints=true,linecolor=gray,linewidth=0.7pt, dotstyle=|, dotangle=90, dotsize=0.15, linestyle=dotted]{%
    1 18.9
2 18.4
3 19.7
4 18.8
5 18.3
}

% Moyenne de la classe
\listplot[plotstyle=line,showpoints=true,linecolor=gray,linewidth=0.7pt, dotstyle=x, dotsize = 0.15, linestyle=dashed]{%
    1 12.4466666666667
2 13.1933333333333
3 10.82
4 11.5066666666667
5 10.1733333333333
}
\end{pspicture}
\pagebreak
\begin{tcolorbox}[enhanced,width=\textwidth,center upper,fontupper=\bfseries,drop shadow southwest,sharp corners]
{\sc \large Hoarau} Alessandro
\end{tcolorbox}
\medskip
\begin{tabularx}{\textwidth}{p{5cm}X}
	\alertbox{\faAward}{Note}{
		\begin{itemize}[leftmargin=0pt]
			\item[\textbullet] Note : \textbf{\large 6.1}
			\item[\textbullet] Rang : \textbf{11}
			\item[\textbullet] Traité : 42 \%
		\end{itemize}
	} &
	\alertbox{\faChartLine}{Statistiques des notes}{
        \psset{xunit=1cm, yunit=1cm,fillstyle=solid}
		\begin{pspicture}(0,-0.1)(16,1.45)
		   \savedata{\data}[7.6 6.6 17.9 4.5 3.5 5.3 16.6 2.7 6.1 11.6 18.0 18.0 18.3 6.3 9.6]
		   \rput{-90}(0,0.9){\psBoxplot[barwidth=1.1cm,yunit=0.5,fillcolor=gray,linewidth=1pt]{\data}}
		   \psaxes[yAxis=false,dx=1cm,Dx=2,labelsep=1pt,linecolor=gray,xlabelFontSize=\scriptstyle](0,0)(10.1,4)
		   \psdot[dotsize=8pt,dotstyle=diamond,linecolor=black,fillstyle=solid,fillcolor=white,linewidth=1pt](3.05,0.85)
           \psdot[dotsize=6pt,dotstyle=x,linecolor=black,linewidth=3pt](5.086666666666667,0.85)
		   \end{pspicture}
	} \\
    
\end{tabularx}\\
\begin{tabularx}{\textwidth}{X}
\alertbox{\faComment}{Commentaire}
{
	Le cours ne semble pas appris, l’exercice 3 n’a presque pas été traité !
}
\end{tabularx}
\medskip
    \ding{113} \textbf{\sffamily{Résultats par thème}} \medskip \\
    \renewcommand{\arraystretch}{1.2}
    \begin{tabular}{|l|r|r|}
    \cline{2-3}
    \multicolumn{1}{l|}{} & \multicolumn{1}{|c|}{Points} & \multicolumn{1}{|c|}{Traitées} \\
    \hline
    {Diviser pour régner} & 22\% \;{\small (08/35)} & 66\% \;{\small (2/3)} \\ \hline {Comprendre un algorithme} & 71\% \;{\small (25/35)} & 75\% \;{\small (3/4)} \\ \hline {Complexité d'un algorithme} & 11\% \;{\small (10/85)} & 12\% \;{\small (1/8)} \\ \hline {Preuve de terminaison} & 0\% \;{\small (00/15)} & 0\% \;{\small (0/1)} \\ \hline {Base d'OCaml (fonctionnel)} & 24\% \;{\small (16/65)} & 33\% \;{\small (2/6)} \\ \hline {Types structurés en C et pointeurs} & 100\% \;{\small (10/10)} & 100\% \;{\small (1/1)} \\ \hline {Programmation de base en C} & 100\% \;{\small (20/20)} & 100\% \;{\small (3/3)} \\ \hline {Programmation dynamique} & 0\% \;{\small (00/40)} & 0\% \;{\small (0/4)} \\ \hline {Ordre et variant} & 43\% \;{\small (13/30)} & 66\% \;{\small (2/3)} \\ \hline \end{tabular} \\\\\medskip \\
    \ding{113} \textbf{\sffamily{Résultats par exercice}} \medskip \\
    \renewcommand{\arraystretch}{1.2}
    \begin{tabular}{|l|r|r|}
    \cline{2-3}
    \multicolumn{1}{l|}{} & \multicolumn{1}{|c|}{Points} & \multicolumn{1}{|c|}{Traitées} \\
    \hline
    Exercice {1} & 43\% \;{\small (13/30)} & 66\% \;{\small (2/3)} \\ \hline Exercice {2} & 50\% \;{\small (15/30)} & 66\% \;{\small (2/3)} \\ \hline Exercice {3} & 5\% \;{\small (06/115)} & 9\% \;{\small (1/11)} \\ \hline Exercice {4} & 42\% \;{\small (68/160)} & 56\% \;{\small (9/16)} \\ \hline \end{tabular} \\\\
    \vspace{0.5cm}\\
    \ding{113} \textbf{\sffamily{Historique des notes}} \medskip \\
    \psset{xunit=1.4cm, yunit=0.2cm}
    \begin{pspicture}(-1,-1)(12,22)

% --- Axe X : numéros de devoir ---
\psaxes[
    Dx=1,
    Dy=2,
    Ox=0,
    Oy=0,
    labels=all,
    ticks=all,
](0,0)(0,0)(7,20)


\listplot[plotstyle=line,showpoints=true,linecolor=black,linewidth=0.7pt, dotstyle=diamond*, dotsize=0.2]{%
    1 8
2 8.2
3 7.5
4 9.6
5 6.1
}

% Minimum de la classe
\listplot[plotstyle=line,showpoints=true,linecolor=gray,linewidth=0.7pt, dotstyle=|, dotangle=90, dotsize=0.15, linestyle=dotted]{%
    1 5.8
2 4.6
3 0.9
4 0
5 2.7
}

% Maximum de la classe
\listplot[plotstyle=line,showpoints=true,linecolor=gray,linewidth=0.7pt, dotstyle=|, dotangle=90, dotsize=0.15, linestyle=dotted]{%
    1 18.9
2 18.4
3 19.7
4 18.8
5 18.3
}

% Moyenne de la classe
\listplot[plotstyle=line,showpoints=true,linecolor=gray,linewidth=0.7pt, dotstyle=x, dotsize = 0.15, linestyle=dashed]{%
    1 12.4466666666667
2 13.1933333333333
3 10.82
4 11.5066666666667
5 10.1733333333333
}
\end{pspicture}
\pagebreak
\begin{tcolorbox}[enhanced,width=\textwidth,center upper,fontupper=\bfseries,drop shadow southwest,sharp corners]
{\sc \large Mahomed Issop} Jérémy
\end{tcolorbox}
\medskip
\begin{tabularx}{\textwidth}{p{5cm}X}
	\alertbox{\faAward}{Note}{
		\begin{itemize}[leftmargin=0pt]
			\item[\textbullet] Note : \textbf{\large 11.6}
			\item[\textbullet] Rang : \textbf{6}
			\item[\textbullet] Traité : 70 \%
		\end{itemize}
	} &
	\alertbox{\faChartLine}{Statistiques des notes}{
        \psset{xunit=1cm, yunit=1cm,fillstyle=solid}
		\begin{pspicture}(0,-0.1)(16,1.45)
		   \savedata{\data}[7.6 6.6 17.9 4.5 3.5 5.3 16.6 2.7 6.1 11.6 18.0 18.0 18.3 6.3 9.6]
		   \rput{-90}(0,0.9){\psBoxplot[barwidth=1.1cm,yunit=0.5,fillcolor=gray,linewidth=1pt]{\data}}
		   \psaxes[yAxis=false,dx=1cm,Dx=2,labelsep=1pt,linecolor=gray,xlabelFontSize=\scriptstyle](0,0)(10.1,4)
		   \psdot[dotsize=8pt,dotstyle=diamond,linecolor=black,fillstyle=solid,fillcolor=white,linewidth=1pt](5.8,0.85)
           \psdot[dotsize=6pt,dotstyle=x,linecolor=black,linewidth=3pt](5.086666666666667,0.85)
		   \end{pspicture}
	} \\
    
\end{tabularx}\\
\begin{tabularx}{\textwidth}{X}
\alertbox{\faComment}{Commentaire}
{
	Bon travail, c’est dommage pour l’exercice 4, où tu n’as pas traité suffisamment de question. Tu semble avoir perdu du temps sur la compréhension de la notion de saut dans un tableau.
}
\end{tabularx}
\medskip
    \ding{113} \textbf{\sffamily{Résultats par thème}} \medskip \\
    \renewcommand{\arraystretch}{1.2}
    \begin{tabular}{|l|r|r|}
    \cline{2-3}
    \multicolumn{1}{l|}{} & \multicolumn{1}{|c|}{Points} & \multicolumn{1}{|c|}{Traitées} \\
    \hline
    {Diviser pour régner} & 0\% \;{\small (00/35)} & 0\% \;{\small (0/3)} \\ \hline {Comprendre un algorithme} & 42\% \;{\small (15/35)} & 100\% \;{\small (4/4)} \\ \hline {Complexité d'un algorithme} & 58\% \;{\small (50/85)} & 62\% \;{\small (5/8)} \\ \hline {Preuve de terminaison} & 80\% \;{\small (12/15)} & 100\% \;{\small (1/1)} \\ \hline {Base d'OCaml (fonctionnel)} & 76\% \;{\small (50/65)} & 83\% \;{\small (5/6)} \\ \hline {Types structurés en C et pointeurs} & 0\% \;{\small (00/10)} & 0\% \;{\small (0/1)} \\ \hline {Programmation de base en C} & 75\% \;{\small (15/20)} & 66\% \;{\small (2/3)} \\ \hline {Programmation dynamique} & 62\% \;{\small (25/40)} & 75\% \;{\small (3/4)} \\ \hline {Ordre et variant} & 93\% \;{\small (28/30)} & 100\% \;{\small (3/3)} \\ \hline \end{tabular} \\\\\medskip \\
    \ding{113} \textbf{\sffamily{Résultats par exercice}} \medskip \\
    \renewcommand{\arraystretch}{1.2}
    \begin{tabular}{|l|r|r|}
    \cline{2-3}
    \multicolumn{1}{l|}{} & \multicolumn{1}{|c|}{Points} & \multicolumn{1}{|c|}{Traitées} \\
    \hline
    Exercice {1} & 93\% \;{\small (28/30)} & 100\% \;{\small (3/3)} \\ \hline Exercice {2} & 90\% \;{\small (27/30)} & 100\% \;{\small (3/3)} \\ \hline Exercice {3} & 69\% \;{\small (80/115)} & 81\% \;{\small (9/11)} \\ \hline Exercice {4} & 37\% \;{\small (60/160)} & 50\% \;{\small (8/16)} \\ \hline \end{tabular} \\\\
    \vspace{0.5cm}\\
    \ding{113} \textbf{\sffamily{Historique des notes}} \medskip \\
    \psset{xunit=1.4cm, yunit=0.2cm}
    \begin{pspicture}(-1,-1)(12,22)

% --- Axe X : numéros de devoir ---
\psaxes[
    Dx=1,
    Dy=2,
    Ox=0,
    Oy=0,
    labels=all,
    ticks=all,
](0,0)(0,0)(7,20)


\listplot[plotstyle=line,showpoints=true,linecolor=black,linewidth=0.7pt, dotstyle=diamond*, dotsize=0.2]{%
    1 13.5
2 15
3 14.8
4 14.2
5 11.6
}

% Minimum de la classe
\listplot[plotstyle=line,showpoints=true,linecolor=gray,linewidth=0.7pt, dotstyle=|, dotangle=90, dotsize=0.15, linestyle=dotted]{%
    1 5.8
2 4.6
3 0.9
4 0
5 2.7
}

% Maximum de la classe
\listplot[plotstyle=line,showpoints=true,linecolor=gray,linewidth=0.7pt, dotstyle=|, dotangle=90, dotsize=0.15, linestyle=dotted]{%
    1 18.9
2 18.4
3 19.7
4 18.8
5 18.3
}

% Moyenne de la classe
\listplot[plotstyle=line,showpoints=true,linecolor=gray,linewidth=0.7pt, dotstyle=x, dotsize = 0.15, linestyle=dashed]{%
    1 12.4466666666667
2 13.1933333333333
3 10.82
4 11.5066666666667
5 10.1733333333333
}
\end{pspicture}
\pagebreak
\begin{tcolorbox}[enhanced,width=\textwidth,center upper,fontupper=\bfseries,drop shadow southwest,sharp corners]
{\sc \large Mamodhoussen} Djavad
\end{tcolorbox}
\medskip
\begin{tabularx}{\textwidth}{p{5cm}X}
	\alertbox{\faAward}{Note}{
		\begin{itemize}[leftmargin=0pt]
			\item[\textbullet] Note : \textbf{\large 18.0}
			\item[\textbullet] Rang : \textbf{2}
			\item[\textbullet] Traité : 94 \%
		\end{itemize}
	} &
	\alertbox{\faChartLine}{Statistiques des notes}{
        \psset{xunit=1cm, yunit=1cm,fillstyle=solid}
		\begin{pspicture}(0,-0.1)(16,1.45)
		   \savedata{\data}[7.6 6.6 17.9 4.5 3.5 5.3 16.6 2.7 6.1 11.6 18.0 18.0 18.3 6.3 9.6]
		   \rput{-90}(0,0.9){\psBoxplot[barwidth=1.1cm,yunit=0.5,fillcolor=gray,linewidth=1pt]{\data}}
		   \psaxes[yAxis=false,dx=1cm,Dx=2,labelsep=1pt,linecolor=gray,xlabelFontSize=\scriptstyle](0,0)(10.1,4)
		   \psdot[dotsize=8pt,dotstyle=diamond,linecolor=black,fillstyle=solid,fillcolor=white,linewidth=1pt](9.0,0.85)
           \psdot[dotsize=6pt,dotstyle=x,linecolor=black,linewidth=3pt](5.086666666666667,0.85)
		   \end{pspicture}
	} \\
    
\end{tabularx}\\
\begin{tabularx}{\textwidth}{X}
\alertbox{\faComment}{Commentaire}
{
	Excellent travail
}
\end{tabularx}
\medskip
    \ding{113} \textbf{\sffamily{Résultats par thème}} \medskip \\
    \renewcommand{\arraystretch}{1.2}
    \begin{tabular}{|l|r|r|}
    \cline{2-3}
    \multicolumn{1}{l|}{} & \multicolumn{1}{|c|}{Points} & \multicolumn{1}{|c|}{Traitées} \\
    \hline
    {Diviser pour régner} & 100\% \;{\small (35/35)} & 100\% \;{\small (3/3)} \\ \hline {Comprendre un algorithme} & 97\% \;{\small (34/35)} & 100\% \;{\small (4/4)} \\ \hline {Complexité d'un algorithme} & 81\% \;{\small (69/85)} & 87\% \;{\small (7/8)} \\ \hline {Preuve de terminaison} & 80\% \;{\small (12/15)} & 100\% \;{\small (1/1)} \\ \hline {Base d'OCaml (fonctionnel)} & 95\% \;{\small (62/65)} & 100\% \;{\small (6/6)} \\ \hline {Types structurés en C et pointeurs} & 100\% \;{\small (10/10)} & 100\% \;{\small (1/1)} \\ \hline {Programmation de base en C} & 100\% \;{\small (20/20)} & 100\% \;{\small (3/3)} \\ \hline {Programmation dynamique} & 75\% \;{\small (30/40)} & 75\% \;{\small (3/4)} \\ \hline {Ordre et variant} & 100\% \;{\small (30/30)} & 100\% \;{\small (3/3)} \\ \hline \end{tabular} \\\\\medskip \\
    \ding{113} \textbf{\sffamily{Résultats par exercice}} \medskip \\
    \renewcommand{\arraystretch}{1.2}
    \begin{tabular}{|l|r|r|}
    \cline{2-3}
    \multicolumn{1}{l|}{} & \multicolumn{1}{|c|}{Points} & \multicolumn{1}{|c|}{Traitées} \\
    \hline
    Exercice {1} & 100\% \;{\small (30/30)} & 100\% \;{\small (3/3)} \\ \hline Exercice {2} & 86\% \;{\small (26/30)} & 100\% \;{\small (3/3)} \\ \hline Exercice {3} & 92\% \;{\small (106/115)} & 100\% \;{\small (11/11)} \\ \hline Exercice {4} & 87\% \;{\small (140/160)} & 87\% \;{\small (14/16)} \\ \hline \end{tabular} \\\\
    \vspace{0.5cm}\\
    \ding{113} \textbf{\sffamily{Historique des notes}} \medskip \\
    \psset{xunit=1.4cm, yunit=0.2cm}
    \begin{pspicture}(-1,-1)(12,22)

% --- Axe X : numéros de devoir ---
\psaxes[
    Dx=1,
    Dy=2,
    Ox=0,
    Oy=0,
    labels=all,
    ticks=all,
](0,0)(0,0)(7,20)


\listplot[plotstyle=line,showpoints=true,linecolor=black,linewidth=0.7pt, dotstyle=diamond*, dotsize=0.2]{%
    1 17.8
2 18
3 16.4
4 17.6
5 18
}

% Minimum de la classe
\listplot[plotstyle=line,showpoints=true,linecolor=gray,linewidth=0.7pt, dotstyle=|, dotangle=90, dotsize=0.15, linestyle=dotted]{%
    1 5.8
2 4.6
3 0.9
4 0
5 2.7
}

% Maximum de la classe
\listplot[plotstyle=line,showpoints=true,linecolor=gray,linewidth=0.7pt, dotstyle=|, dotangle=90, dotsize=0.15, linestyle=dotted]{%
    1 18.9
2 18.4
3 19.7
4 18.8
5 18.3
}

% Moyenne de la classe
\listplot[plotstyle=line,showpoints=true,linecolor=gray,linewidth=0.7pt, dotstyle=x, dotsize = 0.15, linestyle=dashed]{%
    1 12.4466666666667
2 13.1933333333333
3 10.82
4 11.5066666666667
5 10.1733333333333
}
\end{pspicture}
\pagebreak
\begin{tcolorbox}[enhanced,width=\textwidth,center upper,fontupper=\bfseries,drop shadow southwest,sharp corners]
{\sc \large Morel} Lucas
\end{tcolorbox}
\medskip
\begin{tabularx}{\textwidth}{p{5cm}X}
	\alertbox{\faAward}{Note}{
		\begin{itemize}[leftmargin=0pt]
			\item[\textbullet] Note : \textbf{\large 18.0}
			\item[\textbullet] Rang : \textbf{3}
			\item[\textbullet] Traité : 97 \%
		\end{itemize}
	} &
	\alertbox{\faChartLine}{Statistiques des notes}{
        \psset{xunit=1cm, yunit=1cm,fillstyle=solid}
		\begin{pspicture}(0,-0.1)(16,1.45)
		   \savedata{\data}[7.6 6.6 17.9 4.5 3.5 5.3 16.6 2.7 6.1 11.6 18.0 18.0 18.3 6.3 9.6]
		   \rput{-90}(0,0.9){\psBoxplot[barwidth=1.1cm,yunit=0.5,fillcolor=gray,linewidth=1pt]{\data}}
		   \psaxes[yAxis=false,dx=1cm,Dx=2,labelsep=1pt,linecolor=gray,xlabelFontSize=\scriptstyle](0,0)(10.1,4)
		   \psdot[dotsize=8pt,dotstyle=diamond,linecolor=black,fillstyle=solid,fillcolor=white,linewidth=1pt](9.0,0.85)
           \psdot[dotsize=6pt,dotstyle=x,linecolor=black,linewidth=3pt](5.086666666666667,0.85)
		   \end{pspicture}
	} \\
    
\end{tabularx}\\
\begin{tabularx}{\textwidth}{X}
\alertbox{\faComment}{Commentaire}
{
	Excellent travail, l’exercice 3 est vraiment très bien traité, bravo !
}
\end{tabularx}
\medskip
    \ding{113} \textbf{\sffamily{Résultats par thème}} \medskip \\
    \renewcommand{\arraystretch}{1.2}
    \begin{tabular}{|l|r|r|}
    \cline{2-3}
    \multicolumn{1}{l|}{} & \multicolumn{1}{|c|}{Points} & \multicolumn{1}{|c|}{Traitées} \\
    \hline
    {Diviser pour régner} & 100\% \;{\small (35/35)} & 100\% \;{\small (3/3)} \\ \hline {Comprendre un algorithme} & 100\% \;{\small (35/35)} & 100\% \;{\small (4/4)} \\ \hline {Complexité d'un algorithme} & 88\% \;{\small (75/85)} & 87\% \;{\small (7/8)} \\ \hline {Preuve de terminaison} & 80\% \;{\small (12/15)} & 100\% \;{\small (1/1)} \\ \hline {Base d'OCaml (fonctionnel)} & 100\% \;{\small (65/65)} & 100\% \;{\small (6/6)} \\ \hline {Types structurés en C et pointeurs} & 100\% \;{\small (10/10)} & 100\% \;{\small (1/1)} \\ \hline {Programmation de base en C} & 100\% \;{\small (20/20)} & 100\% \;{\small (3/3)} \\ \hline {Programmation dynamique} & 62\% \;{\small (25/40)} & 100\% \;{\small (4/4)} \\ \hline {Ordre et variant} & 80\% \;{\small (24/30)} & 100\% \;{\small (3/3)} \\ \hline \end{tabular} \\\\\medskip \\
    \ding{113} \textbf{\sffamily{Résultats par exercice}} \medskip \\
    \renewcommand{\arraystretch}{1.2}
    \begin{tabular}{|l|r|r|}
    \cline{2-3}
    \multicolumn{1}{l|}{} & \multicolumn{1}{|c|}{Points} & \multicolumn{1}{|c|}{Traitées} \\
    \hline
    Exercice {1} & 80\% \;{\small (24/30)} & 100\% \;{\small (3/3)} \\ \hline Exercice {2} & 90\% \;{\small (27/30)} & 100\% \;{\small (3/3)} \\ \hline Exercice {3} & 100\% \;{\small (115/115)} & 100\% \;{\small (11/11)} \\ \hline Exercice {4} & 84\% \;{\small (135/160)} & 93\% \;{\small (15/16)} \\ \hline \end{tabular} \\\\
    \vspace{0.5cm}\\
    \ding{113} \textbf{\sffamily{Historique des notes}} \medskip \\
    \psset{xunit=1.4cm, yunit=0.2cm}
    \begin{pspicture}(-1,-1)(12,22)

% --- Axe X : numéros de devoir ---
\psaxes[
    Dx=1,
    Dy=2,
    Ox=0,
    Oy=0,
    labels=all,
    ticks=all,
](0,0)(0,0)(7,20)


\listplot[plotstyle=line,showpoints=true,linecolor=black,linewidth=0.7pt, dotstyle=diamond*, dotsize=0.2]{%
    1 16.7
2 15.5
3 13.6
4 14.9
5 18
}

% Minimum de la classe
\listplot[plotstyle=line,showpoints=true,linecolor=gray,linewidth=0.7pt, dotstyle=|, dotangle=90, dotsize=0.15, linestyle=dotted]{%
    1 5.8
2 4.6
3 0.9
4 0
5 2.7
}

% Maximum de la classe
\listplot[plotstyle=line,showpoints=true,linecolor=gray,linewidth=0.7pt, dotstyle=|, dotangle=90, dotsize=0.15, linestyle=dotted]{%
    1 18.9
2 18.4
3 19.7
4 18.8
5 18.3
}

% Moyenne de la classe
\listplot[plotstyle=line,showpoints=true,linecolor=gray,linewidth=0.7pt, dotstyle=x, dotsize = 0.15, linestyle=dashed]{%
    1 12.4466666666667
2 13.1933333333333
3 10.82
4 11.5066666666667
5 10.1733333333333
}
\end{pspicture}
\pagebreak
\begin{tcolorbox}[enhanced,width=\textwidth,center upper,fontupper=\bfseries,drop shadow southwest,sharp corners]
{\sc \large Randriamiarivola Korodo} Lionel
\end{tcolorbox}
\medskip
\begin{tabularx}{\textwidth}{p{5cm}X}
	\alertbox{\faAward}{Note}{
		\begin{itemize}[leftmargin=0pt]
			\item[\textbullet] Note : \textbf{\large 18.3}
			\item[\textbullet] Rang : \textbf{1}
			\item[\textbullet] Traité : 97 \%
		\end{itemize}
	} &
	\alertbox{\faChartLine}{Statistiques des notes}{
        \psset{xunit=1cm, yunit=1cm,fillstyle=solid}
		\begin{pspicture}(0,-0.1)(16,1.45)
		   \savedata{\data}[7.6 6.6 17.9 4.5 3.5 5.3 16.6 2.7 6.1 11.6 18.0 18.0 18.3 6.3 9.6]
		   \rput{-90}(0,0.9){\psBoxplot[barwidth=1.1cm,yunit=0.5,fillcolor=gray,linewidth=1pt]{\data}}
		   \psaxes[yAxis=false,dx=1cm,Dx=2,labelsep=1pt,linecolor=gray,xlabelFontSize=\scriptstyle](0,0)(10.1,4)
		   \psdot[dotsize=8pt,dotstyle=diamond,linecolor=black,fillstyle=solid,fillcolor=white,linewidth=1pt](9.15,0.85)
           \psdot[dotsize=6pt,dotstyle=x,linecolor=black,linewidth=3pt](5.086666666666667,0.85)
		   \end{pspicture}
	} \\
    
\end{tabularx}\\
\begin{tabularx}{\textwidth}{X}
\alertbox{\faComment}{Commentaire}
{
	Excellent travail, bravo Lionel
}
\end{tabularx}
\medskip
    \ding{113} \textbf{\sffamily{Résultats par thème}} \medskip \\
    \renewcommand{\arraystretch}{1.2}
    \begin{tabular}{|l|r|r|}
    \cline{2-3}
    \multicolumn{1}{l|}{} & \multicolumn{1}{|c|}{Points} & \multicolumn{1}{|c|}{Traitées} \\
    \hline
    {Diviser pour régner} & 100\% \;{\small (35/35)} & 100\% \;{\small (3/3)} \\ \hline {Comprendre un algorithme} & 100\% \;{\small (35/35)} & 100\% \;{\small (4/4)} \\ \hline {Complexité d'un algorithme} & 88\% \;{\small (75/85)} & 87\% \;{\small (7/8)} \\ \hline {Preuve de terminaison} & 80\% \;{\small (12/15)} & 100\% \;{\small (1/1)} \\ \hline {Base d'OCaml (fonctionnel)} & 100\% \;{\small (65/65)} & 100\% \;{\small (6/6)} \\ \hline {Types structurés en C et pointeurs} & 100\% \;{\small (10/10)} & 100\% \;{\small (1/1)} \\ \hline {Programmation de base en C} & 100\% \;{\small (20/20)} & 100\% \;{\small (3/3)} \\ \hline {Programmation dynamique} & 62\% \;{\small (25/40)} & 100\% \;{\small (4/4)} \\ \hline {Ordre et variant} & 100\% \;{\small (30/30)} & 100\% \;{\small (3/3)} \\ \hline \end{tabular} \\\\\medskip \\
    \ding{113} \textbf{\sffamily{Résultats par exercice}} \medskip \\
    \renewcommand{\arraystretch}{1.2}
    \begin{tabular}{|l|r|r|}
    \cline{2-3}
    \multicolumn{1}{l|}{} & \multicolumn{1}{|c|}{Points} & \multicolumn{1}{|c|}{Traitées} \\
    \hline
    Exercice {1} & 100\% \;{\small (30/30)} & 100\% \;{\small (3/3)} \\ \hline Exercice {2} & 90\% \;{\small (27/30)} & 100\% \;{\small (3/3)} \\ \hline Exercice {3} & 100\% \;{\small (115/115)} & 100\% \;{\small (11/11)} \\ \hline Exercice {4} & 84\% \;{\small (135/160)} & 93\% \;{\small (15/16)} \\ \hline \end{tabular} \\\\
    \vspace{0.5cm}\\
    \ding{113} \textbf{\sffamily{Historique des notes}} \medskip \\
    \psset{xunit=1.4cm, yunit=0.2cm}
    \begin{pspicture}(-1,-1)(12,22)

% --- Axe X : numéros de devoir ---
\psaxes[
    Dx=1,
    Dy=2,
    Ox=0,
    Oy=0,
    labels=all,
    ticks=all,
](0,0)(0,0)(7,20)


\listplot[plotstyle=line,showpoints=true,linecolor=black,linewidth=0.7pt, dotstyle=diamond*, dotsize=0.2]{%
    1 18.6
2 17.5
3 19.7
4 18.6
5 18.3
}

% Minimum de la classe
\listplot[plotstyle=line,showpoints=true,linecolor=gray,linewidth=0.7pt, dotstyle=|, dotangle=90, dotsize=0.15, linestyle=dotted]{%
    1 5.8
2 4.6
3 0.9
4 0
5 2.7
}

% Maximum de la classe
\listplot[plotstyle=line,showpoints=true,linecolor=gray,linewidth=0.7pt, dotstyle=|, dotangle=90, dotsize=0.15, linestyle=dotted]{%
    1 18.9
2 18.4
3 19.7
4 18.8
5 18.3
}

% Moyenne de la classe
\listplot[plotstyle=line,showpoints=true,linecolor=gray,linewidth=0.7pt, dotstyle=x, dotsize = 0.15, linestyle=dashed]{%
    1 12.4466666666667
2 13.1933333333333
3 10.82
4 11.5066666666667
5 10.1733333333333
}
\end{pspicture}
\pagebreak
\begin{tcolorbox}[enhanced,width=\textwidth,center upper,fontupper=\bfseries,drop shadow southwest,sharp corners]
{\sc \large Rasolofotsara} Ando
\end{tcolorbox}
\medskip
\begin{tabularx}{\textwidth}{p{5cm}X}
	\alertbox{\faAward}{Note}{
		\begin{itemize}[leftmargin=0pt]
			\item[\textbullet] Note : \textbf{\large 6.3}
			\item[\textbullet] Rang : \textbf{10}
			\item[\textbullet] Traité : 52 \%
		\end{itemize}
	} &
	\alertbox{\faChartLine}{Statistiques des notes}{
        \psset{xunit=1cm, yunit=1cm,fillstyle=solid}
		\begin{pspicture}(0,-0.1)(16,1.45)
		   \savedata{\data}[7.6 6.6 17.9 4.5 3.5 5.3 16.6 2.7 6.1 11.6 18.0 18.0 18.3 6.3 9.6]
		   \rput{-90}(0,0.9){\psBoxplot[barwidth=1.1cm,yunit=0.5,fillcolor=gray,linewidth=1pt]{\data}}
		   \psaxes[yAxis=false,dx=1cm,Dx=2,labelsep=1pt,linecolor=gray,xlabelFontSize=\scriptstyle](0,0)(10.1,4)
		   \psdot[dotsize=8pt,dotstyle=diamond,linecolor=black,fillstyle=solid,fillcolor=white,linewidth=1pt](3.15,0.85)
           \psdot[dotsize=6pt,dotstyle=x,linecolor=black,linewidth=3pt](5.086666666666667,0.85)
		   \end{pspicture}
	} \\
    
\end{tabularx}\\
\begin{tabularx}{\textwidth}{X}
\alertbox{\faComment}{Commentaire}
{
	Il faut gagner en rapidité et essayer de traiter plus de questions (notamment celles qui concernent la complexité des algorithmes)
}
\end{tabularx}
\medskip
    \ding{113} \textbf{\sffamily{Résultats par thème}} \medskip \\
    \renewcommand{\arraystretch}{1.2}
    \begin{tabular}{|l|r|r|}
    \cline{2-3}
    \multicolumn{1}{l|}{} & \multicolumn{1}{|c|}{Points} & \multicolumn{1}{|c|}{Traitées} \\
    \hline
    {Diviser pour régner} & 0\% \;{\small (00/35)} & 0\% \;{\small (0/3)} \\ \hline {Comprendre un algorithme} & 48\% \;{\small (17/35)} & 100\% \;{\small (4/4)} \\ \hline {Complexité d'un algorithme} & 18\% \;{\small (16/85)} & 25\% \;{\small (2/8)} \\ \hline {Preuve de terminaison} & 0\% \;{\small (00/15)} & 0\% \;{\small (0/1)} \\ \hline {Base d'OCaml (fonctionnel)} & 61\% \;{\small (40/65)} & 66\% \;{\small (4/6)} \\ \hline {Types structurés en C et pointeurs} & 20\% \;{\small (02/10)} & 100\% \;{\small (1/1)} \\ \hline {Programmation de base en C} & 100\% \;{\small (20/20)} & 100\% \;{\small (3/3)} \\ \hline {Programmation dynamique} & 0\% \;{\small (00/40)} & 0\% \;{\small (0/4)} \\ \hline {Ordre et variant} & 36\% \;{\small (11/30)} & 100\% \;{\small (3/3)} \\ \hline \end{tabular} \\\\\medskip \\
    \ding{113} \textbf{\sffamily{Résultats par exercice}} \medskip \\
    \renewcommand{\arraystretch}{1.2}
    \begin{tabular}{|l|r|r|}
    \cline{2-3}
    \multicolumn{1}{l|}{} & \multicolumn{1}{|c|}{Points} & \multicolumn{1}{|c|}{Traitées} \\
    \hline
    Exercice {1} & 36\% \;{\small (11/30)} & 100\% \;{\small (3/3)} \\ \hline Exercice {2} & 50\% \;{\small (15/30)} & 66\% \;{\small (2/3)} \\ \hline Exercice {3} & 36\% \;{\small (42/115)} & 45\% \;{\small (5/11)} \\ \hline Exercice {4} & 23\% \;{\small (38/160)} & 43\% \;{\small (7/16)} \\ \hline \end{tabular} \\\\
    \vspace{0.5cm}\\
    \ding{113} \textbf{\sffamily{Historique des notes}} \medskip \\
    \psset{xunit=1.4cm, yunit=0.2cm}
    \begin{pspicture}(-1,-1)(12,22)

% --- Axe X : numéros de devoir ---
\psaxes[
    Dx=1,
    Dy=2,
    Ox=0,
    Oy=0,
    labels=all,
    ticks=all,
](0,0)(0,0)(7,20)


\listplot[plotstyle=line,showpoints=true,linecolor=black,linewidth=0.7pt, dotstyle=diamond*, dotsize=0.2]{%
    1 9.5
2 10
3 7
4 6.4
5 6.3
}

% Minimum de la classe
\listplot[plotstyle=line,showpoints=true,linecolor=gray,linewidth=0.7pt, dotstyle=|, dotangle=90, dotsize=0.15, linestyle=dotted]{%
    1 5.8
2 4.6
3 0.9
4 0
5 2.7
}

% Maximum de la classe
\listplot[plotstyle=line,showpoints=true,linecolor=gray,linewidth=0.7pt, dotstyle=|, dotangle=90, dotsize=0.15, linestyle=dotted]{%
    1 18.9
2 18.4
3 19.7
4 18.8
5 18.3
}

% Moyenne de la classe
\listplot[plotstyle=line,showpoints=true,linecolor=gray,linewidth=0.7pt, dotstyle=x, dotsize = 0.15, linestyle=dashed]{%
    1 12.4466666666667
2 13.1933333333333
3 10.82
4 11.5066666666667
5 10.1733333333333
}
\end{pspicture}
\pagebreak
\begin{tcolorbox}[enhanced,width=\textwidth,center upper,fontupper=\bfseries,drop shadow southwest,sharp corners]
{\sc \large Silotia} Donovan
\end{tcolorbox}
\medskip
\begin{tabularx}{\textwidth}{p{5cm}X}
	\alertbox{\faAward}{Note}{
		\begin{itemize}[leftmargin=0pt]
			\item[\textbullet] Note : \textbf{\large 9.6}
			\item[\textbullet] Rang : \textbf{7}
			\item[\textbullet] Traité : 70 \%
		\end{itemize}
	} &
	\alertbox{\faChartLine}{Statistiques des notes}{
        \psset{xunit=1cm, yunit=1cm,fillstyle=solid}
		\begin{pspicture}(0,-0.1)(16,1.45)
		   \savedata{\data}[7.6 6.6 17.9 4.5 3.5 5.3 16.6 2.7 6.1 11.6 18.0 18.0 18.3 6.3 9.6]
		   \rput{-90}(0,0.9){\psBoxplot[barwidth=1.1cm,yunit=0.5,fillcolor=gray,linewidth=1pt]{\data}}
		   \psaxes[yAxis=false,dx=1cm,Dx=2,labelsep=1pt,linecolor=gray,xlabelFontSize=\scriptstyle](0,0)(10.1,4)
		   \psdot[dotsize=8pt,dotstyle=diamond,linecolor=black,fillstyle=solid,fillcolor=white,linewidth=1pt](4.8,0.85)
           \psdot[dotsize=6pt,dotstyle=x,linecolor=black,linewidth=3pt](5.086666666666667,0.85)
		   \end{pspicture}
	} \\
    
\end{tabularx}\\
\begin{tabularx}{\textwidth}{X}
\alertbox{\faComment}{Commentaire}
{
	En programmation c’est satisfaisant. Les aspects mathématiques et le cours doivent être revus.
}
\end{tabularx}
\medskip
    \ding{113} \textbf{\sffamily{Résultats par thème}} \medskip \\
    \renewcommand{\arraystretch}{1.2}
    \begin{tabular}{|l|r|r|}
    \cline{2-3}
    \multicolumn{1}{l|}{} & \multicolumn{1}{|c|}{Points} & \multicolumn{1}{|c|}{Traitées} \\
    \hline
    {Diviser pour régner} & 57\% \;{\small (20/35)} & 66\% \;{\small (2/3)} \\ \hline {Comprendre un algorithme} & 100\% \;{\small (35/35)} & 100\% \;{\small (4/4)} \\ \hline {Complexité d'un algorithme} & 42\% \;{\small (36/85)} & 62\% \;{\small (5/8)} \\ \hline {Preuve de terminaison} & 0\% \;{\small (00/15)} & 100\% \;{\small (1/1)} \\ \hline {Base d'OCaml (fonctionnel)} & 67\% \;{\small (44/65)} & 83\% \;{\small (5/6)} \\ \hline {Types structurés en C et pointeurs} & 100\% \;{\small (10/10)} & 100\% \;{\small (1/1)} \\ \hline {Programmation de base en C} & 75\% \;{\small (15/20)} & 66\% \;{\small (2/3)} \\ \hline {Programmation dynamique} & 0\% \;{\small (00/40)} & 0\% \;{\small (0/4)} \\ \hline {Ordre et variant} & 0\% \;{\small (00/30)} & 100\% \;{\small (3/3)} \\ \hline \end{tabular} \\\\\medskip \\
    \ding{113} \textbf{\sffamily{Résultats par exercice}} \medskip \\
    \renewcommand{\arraystretch}{1.2}
    \begin{tabular}{|l|r|r|}
    \cline{2-3}
    \multicolumn{1}{l|}{} & \multicolumn{1}{|c|}{Points} & \multicolumn{1}{|c|}{Traitées} \\
    \hline
    Exercice {1} & 0\% \;{\small (00/30)} & 100\% \;{\small (3/3)} \\ \hline Exercice {2} & 50\% \;{\small (15/30)} & 100\% \;{\small (3/3)} \\ \hline Exercice {3} & 60\% \;{\small (70/115)} & 81\% \;{\small (9/11)} \\ \hline Exercice {4} & 46\% \;{\small (75/160)} & 50\% \;{\small (8/16)} \\ \hline \end{tabular} \\\\
    \vspace{0.5cm}\\
    \ding{113} \textbf{\sffamily{Historique des notes}} \medskip \\
    \psset{xunit=1.4cm, yunit=0.2cm}
    \begin{pspicture}(-1,-1)(12,22)

% --- Axe X : numéros de devoir ---
\psaxes[
    Dx=1,
    Dy=2,
    Ox=0,
    Oy=0,
    labels=all,
    ticks=all,
](0,0)(0,0)(7,20)


\listplot[plotstyle=line,showpoints=true,linecolor=black,linewidth=0.7pt, dotstyle=diamond*, dotsize=0.2]{%
    1 11.6
2 15.4
3 9
4 0
5 9.6
}

% Minimum de la classe
\listplot[plotstyle=line,showpoints=true,linecolor=gray,linewidth=0.7pt, dotstyle=|, dotangle=90, dotsize=0.15, linestyle=dotted]{%
    1 5.8
2 4.6
3 0.9
4 0
5 2.7
}

% Maximum de la classe
\listplot[plotstyle=line,showpoints=true,linecolor=gray,linewidth=0.7pt, dotstyle=|, dotangle=90, dotsize=0.15, linestyle=dotted]{%
    1 18.9
2 18.4
3 19.7
4 18.8
5 18.3
}

% Moyenne de la classe
\listplot[plotstyle=line,showpoints=true,linecolor=gray,linewidth=0.7pt, dotstyle=x, dotsize = 0.15, linestyle=dashed]{%
    1 12.4466666666667
2 13.1933333333333
3 10.82
4 11.5066666666667
5 10.1733333333333
}
\end{pspicture}
\pagebreak\end{document}